\section{Conclusion}

The pressing question for sample-based quantum computation is no longer whether it can match classical
methods on the ground-state energy---the classical frontier has answered that in the
negative~\cite{leechan2023,belagali2026,reinholdt2025}---but what its output should be, and what
resource actually sets its cost. This work gives a two-part answer. We introduced a method that
computes dynamical spectral functions $A(\omega)$, $A(k,\omega)$, and $S(q,\omega)$ from purely
bitstring-sampled quantum subspaces, lifting the SQD/QSCI paradigm from static energies to
frequency-resolved dynamics with no Hadamard tests, validated against exact diagonalization across
Hubbard chains and nineteen molecules and demonstrated on the IBM~Heron processor---the $L=6$ Hubbard
spectral function $A(\omega)$ reconstructed on \texttt{ibm\_fez}, and, for the ground-state energy of
stretched $\mathrm{N_2}$, a noise-assisted configuration-recovery run on \texttt{ibm\_marrakesh}
reaching $0.59\pm0.14$~mHa of the exact value (over $8$ recovery seeds).

We then asked which resource controls the determinant support $|\mathcal S|$ that the sampler must
populate, and corrected a natural expectation. Restricted to the number-conserving states of
chemistry, the fermionic AntiFlatness collapses to a single one-particle-density-matrix invariant,
$\mathcal F_1=4\,\mathrm{tr}[\gamma(1-\gamma)]$ (equal to $2N_{\mathrm u}$ for spin-restricted natural
orbitals), identifying fermionic magic with the effectively-unpaired-electron index and one-body entanglement. Being a Gaussian (orbital-rotation)
invariant while $|\mathcal S|$ is basis dependent, it is provably decoupled from the sampling
cost---two-sidedly and without a stable sign---and no orbital-rotation-invariant functional of the
reduced density matrices can repair this. Across the nineteen-molecule suite one-body magic is nonetheless a faithful,
efficiently measurable detector of static, multireference character---switching on for covalent
bond-breaking, already large for inherently multireference $\mathrm{C_2}$ and $\mathrm{O_2}$,
near-Gaussian for the ionic bond of LiF---but it is a diagnostic of correlation, not a predictor of
classical cost---which is instead lower-bounded and, on the systems realized here, empirically tracked by the basis-dependent entanglement, the bond dimension $\chi$.
Any genuine quantum advantage lives one stratum above, in the generic, high-rank non-Gaussianity of the
connected higher-body cumulants in the fixed sampling basis---not any orbital-invariant magic order---beyond
the stratum whose free-fermionic estimation primitives remain classically tractable~\cite{oh2026}.

We have been deliberate about scope. The reach of the method is asymptotic and lives in two-dimensional
and chemical settings rather than one dimension, where tensor networks remain near-optimal~\cite{schollwock2011}; it rests on
a moment-exactness and polynomial-shot theorem (Thm.~\ref{thm:b2}, App.~\ref{app:b2}); and, on the artificial-intelligence side of the hybrid
loop, self-consistent configuration recovery improves the sampled subspace under realistic device
noise through valid-shot economy rather than quantum advantage, the large effect in the dequantization
test is a free classical one---an Epstein--Nesbet tail on the existing device histogram, no extra shots,
cutting the energy error by $31\%$ ($39.24\pm3.95\to26.98\pm0.66$~mHa)---and a \emph{learned} generative
tail-completer does \emph{not} beat that corrected baseline (Fig.~\ref{fig:gflow})---delimiting, not
foreclosing, the quantum--AI frontier, and consistent with the classical-simulability boundary. Read together, the two contributions deliver a dynamical primitive the
energy-only hardware of today does not provide, and a precise account of which fermionic resource sets
its cost and which merely diagnoses correlation---handing the quantum-centric supercomputing programs
now operating at industrial scale~\cite{robledomoreno2025,shirakawa2025} a computable multireference
diagnostic---and a proof that this easily-measured quantity does not by itself map the classical cost
frontier, which is instead lower-bounded and empirically tracked by the bond dimension $\chi$. The natural next steps are to push the sampled-subspace reconstruction into
the two-dimensional and multi-orbital regimes where the higher-order hardness lives, and to make the
\emph{basis-dependent} cost---the bond dimension $\chi$ in the sampler's own frame, and the generic
high-rank non-Gaussianity above it---not the orbital-invariant one-body magic that is so easily
measured---the object one estimates when deciding whether a quantum sampler is worth its shots.
